\section{Kết nối hồi quy và phân lớp}
\subsection{Logistic Regression như một bài toán hồi quy}
Mặc dù thường được sử dụng cho bài toán phân loại nhị phân, Logistic
Regression thực chất là một dạng hồi quy phi tuyến. Cụ thể, mô hình giả định
rằng xác suất của biến mục tiêu $t \in \{0,1\}$ được xác định thông qua một
hàm sigmoid áp dụng lên tổ hợp tuyến tính của đầu vào:
\begin{equation}
p(t=1|\boldsymbol{x}) = \sigma(\boldsymbol{w}^\intercal \boldsymbol{x}) =
\frac {1}{1 + e^{-\boldsymbol{w}^\intercal \boldsymbol{x}}}.
\end{equation}
Do đó, thay vì dự đoán trực tiếp nhãn lớp, mô hình thực hiện hồi quy trên xác
suất (hoặc kỳ vọng) của $t$. Điều này cho thấy Logistic Regression là một mô
hình hồi quy với hàm liên kết phi tuyến (sigmoid), ánh xạ từ không gian tuyến
tính sang miền xác suất $[0,1]$.

\subsection{Khung Generalized Linear Models (GLM)}
Generalized Linear Models (GLM) cung cấp một khung thống nhất cho nhiều mô
hình hồi quy khác nhau \cite[Chap.~12]{pml1Book}. Trong GLM, ta giả định rằng:
\begin{itemize}
    \item Biến đầu ra $t$ tuân theo một phân phối thuộc exponential family.
    \item Kỳ vọng của $t$, ký hiệu $\mu = \mathbb{E}[t|\boldsymbol{x}]$, liên
    hệ với tổ hợp tuyến tính $\eta = \boldsymbol{w}^\intercal \boldsymbol{x}$
    thông qua một hàm liên kết $g$:
    \begin{equation}
    g(\mu) = \eta = \boldsymbol{w}^\intercal \boldsymbol{x}.
    \end{equation}
\end{itemize}
Các mô hình quen thuộc là các trường hợp đặc biệt của GLM:
\begin{itemize}
    \item Linear Regression: phân phối Gaussian, link function là hàm đồng
    nhất $g(\mu) = \mu$.
    \item Logistic Regression: phân phối Bernoulli, link function là logit $g
    (\mu) = \log \frac{\mu}{1-\mu}$.
    \item Poisson Regression: phân phối Poisson, link function là log $g(\mu)
    = \log \mu$.
\end{itemize}
Như vậy, GLM cho phép ta thay đổi phân phối của dữ liệu đầu ra và hàm liên kết
một cách linh hoạt nhưng vẫn giữ cấu trúc tuyến tính trong không gian tham số.

\subsection{Vai trò của exponential family}
Một điểm then chốt của GLM là giả định rằng phân phối của $t$ thuộc
exponential family, có dạng:
\begin{equation}
p(t|\theta) = h(t)\exp\left(\theta^\intercal \boldsymbol{\phi}(t) - A(\theta)
\right).
\end{equation}
Các phân phối như Gaussian, Bernoulli, và Poisson đều thuộc họ này. Điều này
mang lại một cấu trúc toán học thống nhất với các lợi ích sau:
\begin{itemize}
    \item Kỳ vọng và phương sai có thể được suy ra trực tiếp từ hàm $A(\theta)
    $:
    \begin{equation}
    \mathbb{E}[t] = A'(\theta), \quad \mathrm{Var}(t) = A''(\theta).
    \end{equation}
    \item Hàm log-likelihood có dạng lồi đối với nhiều trường hợp, giúp tối ưu
    hóa hiệu quả.
    \item Tồn tại các sufficient statistics hữu hạn chiều, giúp nén dữ liệu mà
    không mất thông tin.
\end{itemize}
Do đó, cả Linear Regression (Gaussian) và Logistic Regression (Bernoulli) đều
nằm trong một khuôn khổ thống nhất, nơi sự khác biệt chỉ nằm ở lựa chọn phân
phối và hàm liên kết.

\subsection{Tổng kết}
Logistic Regression có thể được hiểu là một dạng hồi quy với hàm sigmoid. GLM
mở rộng ý tưởng này bằng cách cho phép nhiều phân phối đầu ra khác nhau thông
qua exponential family và các hàm liên kết tương ứng. Chính exponential family
cung cấp nền tảng lý thuyết giúp hợp nhất các mô hình này thành một hệ thống
chung, vừa linh hoạt vừa có cấu trúc toán học chặt chẽ.
